\section{MPI, parallel parameters and sharing a GPU}
\label{sec:mpi}

\subsection{What MPI does in a GPU run}

ALF's MPI parallelism is over independent Markov chains: each rank runs its own sweep, and only the observables are reduced at the end of a bin. The GPU port keeps this. Every rank builds its own device mirror and its own device $G$ and stacks, draws its own random numbers on the host, and talks to the other ranks only through the host-side observable reduction (\file{gpu\_state\_mod.F90:516-519}). No GPU-aware MPI is needed.

\paragraph{Rank-to-device binding.} The code sets the OpenMP default device to
\begin{lstlisting}
call omp_set_default_device(mod(env_rank, omp_get_num_devices()))
\end{lstlisting}
(\file{gpu\_state\_mod.F90:550-552}), where \texttt{env\_rank} is the rank in \texttt{MPI\_COMM\_WORLD}. In practice the job scripts restrict each rank to one device before the program starts, so the rule always yields device 0:
\begin{itemize}
\item NVIDIA (Helma): a wrapper sets \env{CUDA\_VISIBLE\_DEVICES} $=r/\text{ranks per GPU}$ from the Open MPI rank $r$, and with MPS the pipe and log directory of that GPU's MPS daemon (\file{docs/gpu/scripts/helma/rank\_wrap.sh}). The device used by the OpenACC stream query (Section~\ref{sec:design}) must be the OpenMP default device, which this binding guarantees.
\item Intel (SuperMUC-NG Phase 2): \env{ZE\_FLAT\_DEVICE\_HIERARCHY=FLAT} exposes the 8 tiles of a node's four Max~1550 cards as 8 devices, so rank $r$ uses tile $r\bmod8$; with fewer ranks \env{ZE\_AFFINITY\_MASK} binds them (\file{intel\_pvc\_supermuc\_ng2.md:26-27}).
\item AMD: no MPI run is documented (\file{FEATURES.md}, line 47).
\end{itemize}

\dis{The binding uses the global rank, not the node-local rank. On several nodes the mapping is right only if the ranks are placed in blocks and the number of ranks per node is a multiple of the number of devices. This is not documented; the job scripts avoid the issue by setting the visible devices per rank.}

\paragraph{Correctness.} One MPI rank reproduces the non-MPI run to $10^{-12}$ (energy $7\times10^{-12}$, \texttt{SpinZ\_localtau} $1.4\times10^{-12}$, \texttt{SpinZ\_local} $7\times10^{-15}$), and acceptance (0.3375) and precision ($10^{-14}$) do not depend on the number of ranks (\file{MPI\_GPU.md}). % MPI_GPU.md

\subsection{Parallel parameters and tempering}

With \texttt{TEMPERING} and \texttt{PARALLEL\_PARAMS} ALF runs groups of ranks with different parameter sets. Real parallel tempering (exchange of configurations between parameter sets) is refused on the GPU path. Parameter-parallel runs without exchange are allowed: the groups are independent, and \file{main.F90} skips \texttt{Exchange\_Step} (\file{PP\_GPU.md}). Build: \texttt{configure.sh PGI PP GPU=nvidia HDF5} with \texttt{PP = -DMPI -DTEMPERING -DPARALLEL\_PARAMS}.

\subsection{Several ranks on one GPU}

One chain does not fill a GPU: on the A30 one chain kept the GPU busy 58\,\% of the time, two chains 82\,\% (\file{MPI\_GPU.md}, \texttt{nvidia-smi dmon}); on the H100 at $L=32$ one chain used 55\,\% of the streaming multiprocessors (\file{helma\_h100.md}, section 4.2). % MPI_GPU.md; helma_h100.md 4.2
The reason is the structure of the work: most kernels are small (block gathers, rank-$k$ flushes, single-vertex wraps) and the host spends time between them. Running several chains per GPU fills the gaps. How well this works depends on how the GPU shares itself between processes.

\paragraph{NVIDIA: MPS.} Without help, processes on one NVIDIA GPU are \emph{time-sliced}: the GPU switches between their contexts, and kernels of different processes never run at the same time. The \emph{Multi-Process Service} (MPS) merges the processes into one context, so their kernels run concurrently. \env{CUDA\_MPS\_ACTIVE\_THREAD\_PERCENTAGE} limits the share of the GPU each client can use. Measured:
\begin{itemize}
\item Without MPS, two processes on the A30 ran 13 to 15 times slower each (\file{CLAUDE.md}); on the H100 two ranks were 6.5 times slower per chain at $L=24$ (29.5 against 4.54~s per sweep) and 2.6 times at $L=32$; four or eight ranks at $L=32$ did not finish six sweeps in 420~s (\file{helma\_h100.md}, section 4). % helma_h100.md
\item With MPS the throughput rises with the number of ranks. Table~\ref{tab:mps} gives the A30 at $L=16$.
\item A 25\,\% thread share was 12\,\% faster than no limit with 9 ranks at $L=24$ (86.9 against 98.3~s for 10 sweeps, \file{PP\_GPU.md}); 25\,\% and 50\,\% agree within 2\,\% on the H100. % PP_GPU.md; helma_h100.md
\end{itemize}

\begin{table}[h]
\centering\small
\caption{A30, $L=16$, submatrix $\Nd=32$, 30 sweeps, with MPS (\file{MPI\_GPU.md}).}
\label{tab:mps}
\begin{tabular}{@{}rlrrr@{}}
\toprule
Ranks & MPS share & Time per chain (s) & vs 1 rank & Sweeps/s per GPU \\
\midrule
1 & 100\,\% & 38.8 & -- & 0.73 \\ % MPI_GPU.md
2 & 50\,\% & 45.3 & +17\,\% & 1.25 \\
4 & 25\,\% & 54.6 & +41\,\% & 2.09 \\
8 & 13\,\% & 81.4 & +110\,\% & 2.83 \\
16 & default & 145.6 & +275\,\% & 3.18 \\
16 & 25\,\% & -- & -- & 4.68 \\
\bottomrule
\end{tabular}
\end{table}

On the H100 at $L=24$ one rank gives 0.285 sweeps per second, eight ranks with MPS at 25\,\% give 1.176, 4.1 times more, while each chain becomes 1.9 times slower (6.80 against 3.51~s per sweep). At $L=32$: 0.102 with one rank, 0.232 with eight and 0.250 with sixteen (\file{helma\_h100.md}, sections 4.1-4.2). % helma_h100.md
MPS does not change the results: two MPS clients produced output bit-identical to a solo run (\file{CLAUDE.md}).

\paragraph{Intel: CCS.} The Max~1550 has no MPS. Ranks sharing a tile are time-sliced by default. The compute command streamer setting \env{ZEX\_NUMBER\_OF\_CCS=0:4,\ldots,7:4} splits each tile into four hardware queues, so four processes can run concurrently on one tile. With the V7 switches at $L=24$ the throughput per card is 0.365, 0.458, 0.879 and 0.261 sweeps per second for 1, 2, 4 and 8 walkers per tile (\file{ISSUE\_INTEL\_VS\_A30.md}, section 7). % ISSUE section 7
Four walkers per tile (32 ranks per node) is the optimum.

\dis{\file{intel\_pvc\_supermuc\_ng2.md} (line 45, 2026-11-10) reports that CCS 4 gave ``no gain'' and recommends one rank per tile (line 47); \file{RUNTIME\_SWITCHES.md} (line 173) and the later session notes report CCS 4 with 32 ranks per node as the throughput optimum. The later measurement is the one used in production.}

\paragraph{AMD.} Four concurrent processes on the MI210 without any sharing service gave 1.95 times the throughput of one (\file{SPEEDUPS.md}, section 3); compute-unit masking (\env{HSA\_CU\_MASK}) gave nothing. % SPEEDUPS.md

\subsection{The production layout}

The production layout is a measured per-platform setting (\file{GPU\_PLAN\_V2.md}, section 2): on the H100, eight ranks per GPU with MPS at 25\,\%, $\Nd=64$ at $L=32$ (\file{helma\_h100.md}, section 7); on the Max~1550, four ranks per tile with CCS 4; on the A30 the owner decided on one rank per GPU (2026-11-08, \file{MPI\_GPU.md}). The choice of $\Nd$ interacts with the layout: with eight ranks on the H100 at $L=32$, $\Nd=32$ gave 66.5~s per sweep per chain and $\Nd=64$ gave 34.5~s (\file{helma\_h100.md}, section 4.2). % helma_h100.md
